\section*{Author Section 4: resilience under erasures and errors}
For parameters $m,r$, recall that $E(m,r)$ is the $\binom m{\le r}\times2^m$-matrix where the columns are indexed by elements in $\mathbb F_2^m$, and the $z$-th column is the column vector consisting of evaluations of all multilinear $m$-variate degree $\le r$ monomials on $z$. It is well-known that the matrix $E(m,r)$ is the generator matrix of the $\mathrm{RM}(m,r)$ code, and is also the parity-check matrix of the $\mathrm{RM}(m,m-r-1)$ code. Theorem 1.1 can be re-interpreted as a statement about random sub-matrices of $E(m,r)$ in the regime where $r\le\gamma m$ for a small absolute constant $\gamma>0$.
\paragraph{Corollary 4.1 (rank of random sub-matrices of $E(m,r)$).} Let $\gamma>0$ be an absolute constant and $m,r$ growing parameters with $r\le\gamma m$. Then, for any constant $\epsilon>0$, a random set of $K=(1-\epsilon)\binom m{\le r}$ columns of $E(m,r)$ are linearly independent with probability $1-o(1)$.
\begin{proof}A set of columns of $E(m,r)$ indexed by $Z=\{z_1,\ldots,z_K\}$ are linearly dependent if and only if $I_r(Z)$ has dimension strictly larger than $\binom m{\le r}-K$. Theorem 1.1 asserts that this happens with only $o(1)$ probability.\end{proof}
\subsection*{Author Section 4.1: channels under consideration}
The above corollary can be used to talk about the resilience of Reed-Muller codes under the erasure and error channels. We first define the precise definitions to be able to state the results accurately.
\paragraph{Definition 4.2 (binary erasure channel (BEC)).} The binary erasure channel with parameter $\alpha$, denoted by $\mathrm{BEC}_\alpha$, is the channel with input alphabet $\{0,1\}$ where the each binary symbol is ``erased'' (replaced by the `?' symbol) independently with probability $\alpha$.
A closely related model is where we have a fixed number of random erasures instead of each coordinate being erased with a fixed probability. We will refer to this as the $\mathrm{BEC}^*$ model although this isn't a channel in the traditional sense of altering each coordinate independently.
\paragraph{Definition 4.3 (capped binary erasure channel (BEC*)).} The capped binary erasure channel with parameter $K$, denoted by $\mathrm{BEC}_K^*$, refers to the channel with input alphabet $\{0,1\}^n$ that replaces a random subset of at most $K$ of the coordinates by the `?' symbol.
The Binary Symmetric Channel deals with errors or corruptions as opposed to erasures in the Binary Erasure Channel.
\paragraph{Definition 4.4 (binary symmetric channel (BSC)).} The binary symmetric channel with parameter $\alpha$, is the channel where the each binary symbol is ``flipped'' (that is, 0 changed to 1 and vice-versa) independently with probability $\alpha$.
In similar spirit to Definition 4.3, we define the capped binary symmetric channel.
\paragraph{Definition 4.5 (capped binary symmetric channel (BSC*)).} The capped binary symmetric channel with parameter $K$, denoted by $\mathrm{BSC}_K^*$, refers to the channel with input alphabet $\{0,1\}^n$ that ``flips'' a random subset of at most $K$ of the coordinates.
In most cases, resilience with respect to $\mathrm{BEC}_\alpha$ (or $\mathrm{BSC}_\alpha$) is the same as resilience with respect to $\mathrm{BEC}_K^*$ (or $\mathrm{BSC}_K^*$) for $K=\alpha n$, by standard concentration inequalities but this might require some subtlety when dealing with codes of rate very close to zero or close to one. For a concrete example, the code $\mathrm{RM}(m,m-1)$ has rate $R=1-1/2^m$ and is not resilient under $\mathrm{BEC}_\alpha$ for $\alpha=(1-R)/2=1/2^{m+1}$ (with constant probability we may have two coordinates erased, and this is not recoverable), but is vacuously resilient under $\mathrm{BEC}_K^*$ for $K=\alpha2^m$. To avoid these nuances, we will only be dealing with the setting of capped channels although this does not make much difference with most of our range of parameters.
\paragraph{Notion of ``capacity achieving'' codes with respect to $\mathrm{BEC}_\alpha$.} Shannon's seminar work [14] showed that, for any constant rate, the supremum of rates of random codes that is resilient to $\mathrm{BEC}_\alpha$ is exactly $R=1-\alpha$. That is, for any $\epsilon>0$, there are codes of rate $1-\alpha-\epsilon$ that can decode from $\mathrm{BEC}_\alpha$ with decoding error $1-o(1)$. However, when the rate is very close to zero or one, the situation becomes more nuanced due to the asymptotics involved.
Dealing specifically with Reed-Muller codes, let us consider the code $\mathrm{RM}(m,m-r-1)$ with $r\le m/2$. This has rate $R=1-\binom m{\le r}/2^m$ and can recover from a maximum of $\binom m{\le r}$ errors. Abbe, Shpilka and Wigderson [1] defined the notion of ``capacity achieving under BEC'' to mean that the code can decode (with $1-o(1)$ probability) from $(1-\epsilon)\binom m{\le r}$ erasures for any constant $\epsilon>0$. We refer the reader to [1, Section 1] and [13, Section 2.2] for a nuanced discussion on this.
\subsection*{Author Section 4.2: Reed-Muller codes under BEC*}
An immediate consequence of the above corollary is that $\mathrm{RM}(m,d)$ with $d\ge m(1-\gamma)$ for a small but absolute constant $\gamma>0$ achieves capacity (as defined above) under the Binary Erasure Channel.
\paragraph{Proof of Theorem 1.2.} The code $\mathrm{RM}(m,d)$ can recover from erasures on coordinates indexed by $Z\subseteq\mathbb F_2^m$ if and only if the corresponding columns of the parity check matrix are linearly independent. Since the parity-check matrix of $\mathrm{RM}(m,d)$ is $E(m,r)$ for $r=m-d-1$, we have from Corollary 4.1 that a random set of $K$ columns are linearly independent with probability $1-o(1)$.\qed
\subsection*{Author Section 4.3: Reed-Muller codes under BSC*}
The following theorem of Abbe, Shpilka and Wigderson [1] provides a method to derive the resilience of certain Reed-Muller codes under the BSC by using the resilience of appropriate Reed-Muller codes under the BEC. Subsequent work of Saptharishi, Shpilka and Volk also showed that these codes also have efficient decoding procedures to recover from random errors.
\paragraph{Theorem 4.6 ([1,12]).} For any growing parameters $m,r>0$, if the code $\mathrm{RM}(m,m-2r-2)$ can recover from a subset $Z\subseteq[2^m]$ of erasures, then the code $\mathrm{RM}(m,m-r)$ can efficiently recover from errors on the subset $Z$. In particular, if the code $\mathrm{RM}(m,m-2r-2)$ can recover $K$ random erasures with probability $1-o(1)$, then the code $\mathrm{RM}(m,m-r)$ can efficiently decode from $K$ random errors.
Applying the above theorem to Theorem 1.2 yields Corollary 1.3.
\paragraph{Corollary 4.7.} There exists a constant $\gamma_0>0$ such that for all $\gamma<\gamma_0$ and $\epsilon>0$ the following is true. Let $m,r$ be growing parameters with $r=\gamma m$. Then, the code $\mathrm{RM}(m,m-2r-2)$ can be efficiently decode from $K=(1-\epsilon)\binom m{\le r}$ random errors.
It is worth noting that the minimum distance of the code $\mathrm{RM}(m,m-2r-2)$ is merely $2^{2r+2}\ll\binom m{\le r}$ when $r=\gamma m$ for a small enough $\gamma>0$. Thus, the above corollary shows that high-degree (or high-rate) Reed-Muller codes are resilient to random errors well beyond their minimum distance, and efficiently so.
